\documentclass{article}

\usepackage[urw-garamond]{mathdesign}

\usepackage{geometry}
\usepackage{verbatim}
\usepackage{cite}
\usepackage{listings}

\lstdefinelanguage{siexml}
  {morekeywords={xml,sie,version,xmlns,decoder,loop,var,bits,type,endian,
                 value,group,tag,id,format,read,sample,octets,encoding,
                 ch,test,index,dim,set,data,name},
   sensitive=false,
   morecomment=[s]{<!--}{-->},
   morestring=[b]",
   showstringspaces=false,
   columns=fullflexible
  }

\lstdefinelanguage{rng}{
  morekeywords={attribute,div,element,
    empty,external,grammar,include,inherit,list,
    mixed,notAllowed,parent,string,
    text,token},
  sensitive=false,
  morecomment=[l]{\#},
  morestring=[b]",
  columns=fullflexible
}

\lstnewenvironment{grammar}
  {\vbox\bgroup\lstset{language=rng,frame=simple}}
  {\egroup}

\lstnewenvironment{siexmlexample}
  {\vbox\bgroup\lstset{language=siexml,frame=simple}}
  {\egroup}

\newcommand{\relaxng}[0]{\hbox{RELAX~NG}}
\newcommand{\xml}[1]{\texttt{#1}}
\newcommand{\var}[1]{\texttt{#1}}
\newcommand{\lit}[1]{\texttt{#1}}

\newcommand{\fixme}{}
\newcommand{\incomplete}{\emph{(\ldots incomplete\ldots)}}

\newcommand{\undef}{\emph{undefined}}
\newcommand{\libsie}{\xml{libsie}}


\newcommand{\somatver}{1.0}

\title{The \xml{somat} SIE Schema, Version \somatver{}}

\author{
Brian Downing \\ HBM, Inc. \and
Chris LaReau \\ HBM, Inc.
}

\date{December, 2009}

\begin{document}

\maketitle

\section{Introduction}

This document describes the \xml{somat} schema for the SIE format.
This schema is used by SoMat's line of data acquisition system.  In
addition, some elements of this schema may be applicable for more
general use.  In almost all cases data files that reference this
schema will also reference the \xml{core} schema.

\section{Metadata schema}

\subsection{\xml{somat:data\_bits} tag}

This tag contains the number of bits of resolution of the channel data.

\begin{siexmlexample}
<tag id="somat:data_bits">16</tag>
\end{siexmlexample}

\subsection{\xml{somat:data\_format} tag}

This tag contains the source data format of the channel data:
``int'', ``uint'', or ``float''.

\begin{siexmlexample}
<tag id="somat:data_format">uint</tag>
\end{siexmlexample}

\subsection{\xml{somat:datamode\_name} tag}

This tag contains the name of the datamode which produced the current
SIE channel.

\begin{siexmlexample}
<tag id="somat:datamode_name">th1k</tag>
\end{siexmlexample}

\subsection{\xml{somat:datamode\_type} tag}

This tag contains the datamode type of the datamode that produced the
current SIE channel.  It can be one of the following values: \\
  \xml{time\_history} \\
  \xml{burst\_history} \\
  \xml{peak\_valley} \\
  \xml{peak\_valley\_slice} \\
  \xml{event\_slice} \\
  \xml{time\_at\_level} \\
  \xml{peak\_valley\_matrix/from\_to} \\
  \xml{peak\_valley\_matrix/range\_mean} \\
  \xml{peak\_valley\_matrix/range\_only} \\
  \xml{rainflow/from\_to} \\
  \xml{rainflow/range\_mean} \\
  \xml{rainglow/range\_only} \\
  \xml{message\_log}

\begin{siexmlexample}
<tag id="somat:datamode_type">time_history</tag>
\end{siexmlexample}

\subsection{\xml{somat:input\_channel} tag}

This tag contains the name of the input channel which produced the
current SIE channel.

\begin{siexmlexample}
<tag id="somat:input_channel">bracket</tag>
\end{siexmlexample}

\subsection{\xml{somat:log} tag}

This tag contains the data acquisition system's log.  On the eDAQ,
this is the same as what is retrieved with the ``FCS Setup | Upload
FCS Log'' command in TCE.  The log is stored continuously and is as
such not limited in size.

\subsection{\xml{somat:rainflow\_unclosed\_cycles} tag}

This tag contains a sequence of ASCII-formatted floating-point numbers
containing the unclosed cycles stack from the rainflow counting
algorithm.

\begin{siexmlexample}
<tag id="somat:rainflow_unclosed_cycles">
  9.992119789123535 20.02879905700684 -40.09270095825195
  40.0364990234375 -50.09659957885742 60.07699966430664
  -70.13710021972656 70.08080291748047 -80.14089965820312
  90.12129974365234 -110.1849975585938 150.2100067138672
  -120.1890029907227 190.2579956054688 -150.2330017089844
  230.3390045166016 -330.5 390.5650024414062 -681.0609741210938
  630.9199829101562 -590.89501953125 460.6900024414062
  -460.7139892578125 340.5130004882812 -320.4960021972656
  170.2510070800781 -210.3220062255859
</tag>
\end{siexmlexample}

\subsection{\xml{somat:tce\_setup} tag}

This tag contains the .TCE setup file used to initialize the current
test.

%% \subsection{\xml{somat:transducer\_rezero} tag}

%% \incomplete{}

\subsection{\xml{somat:version} tag}

The value of the \xml{somat:version} tag defines the version of the
\xml{somat} schema in use.  The version described in this document is
\somatver{}.

\begin{siexmlexample}
<tag id="somat:version">1.0</tag>
\end{siexmlexample}

\newcommand{\nondecreasing}[1]{The values of SIE dimension index #1,
  scaled or unscaled, have non-decreasing ordering}
\newcommand{\schemaheader} {\textbf{SIE dimension index} & 
  \textbf{Data type} & \textbf{Scaled} & \textbf{Unscaled}}

\section{Data schemas}

\subsection{\xml{somat:sequential} data schema}

The \xml{somat:sequential} data schema is used for representing time
series numerical data sampled at a regular interval, or the sequential
output of various data reduction algorithms which dispose of time
information.  Each row represents a single data sample and the time or
sequence number of that sample.

Where time is preserved, the data output represents:

\begin{center}
  \begin{tabular}{llll}
    \schemaheader{} \\ \hline
    0 & Numeric & Time & Sample number \\
    1 & Numeric & Engineering data value & \undef{}
  \end{tabular}
\end{center}

Otherwise, the data represents:

\begin{center}
  \begin{tabular}{llll}
    \schemaheader{} \\ \hline
    0 & Numeric & Sequence number & Sequence number \\
    1 & Numeric & Engineering data value & \undef{}
  \end{tabular}
\end{center}

\nondecreasing{0}.

\subsection{\xml{somat:message} data schema}

The \xml{somat:message} data schema is used for representing
non-numerical data sampled at irregular intervals.  Each row
represents a single message or event and the time of collection.

\begin{center}
  \begin{tabular}{llll}
    \schemaheader{} \\ \hline
    0 & Numeric & Time & \undef{} \\
    1 & Raw & Binary message data & \undef{}
  \end{tabular}
\end{center}

\nondecreasing{0}.

\subsection{\xml{somat:burst} data schema}

The \xml{somat:burst} data schema is used for representing time series
numerical data sampled at a regular interval, but where actual data
collection is triggered by a triggering event.  Each row represents a
single data sample, the time that sample was collected, and the
relation between that sample and the event which triggered collection.

\begin{center}
  \begin{tabular}{llll}
    \schemaheader{} \\ \hline
    0 & Numeric & Time & Sample number \\
    1 & Numeric & Engineering data value & \undef{} \\
    2 & Numeric & Burst index & Burst index
  \end{tabular}
\end{center}

\nondecreasing{0}.

For \xml{somat:burst}, the \emph{sample number} is absolute; i.e., if
the first burst data happened at the 50,000th sample collected, the
sample number emitted from SIE dimension 0 for that data would be
50,000, not 0.

The purpose of the \emph{burst index} is so that the position of
the burst trigger can be known:

\begin{center}
  \begin{tabular}{ll}
    \textbf{Burst index} & \textbf{Meaning} \\ \hline
    Integer $n < 0$ & Current sample is $n$ samples before 
      the burst trigger. \\
    0 & Current sample is the first sample in the burst trigger. \\
    0.5 & Current sample is a continuation of the (level-sensitive) 
      burst trigger. \\
    Integer $n > 0$ & Current sample is $n$ samples after the 
      burst trigger.
  \end{tabular}
\end{center}

\subsection{\xml{somat:histogram} data schema}

The \xml{somat:histogram} schema is used for representing
$n$-dimensional histogram data.  Each row represents a single
histogram bin's count and limits in all incoming dimensions.  The
histogram bins are not presented in any particular order, nor are
empty bins guaranteed to be present.

\begin{center}
  \begin{tabular}{llll}
    \schemaheader{} \\ \hline
    0 & Numeric & Bin count & \undef{} \\
    $n * 2 + 1$ & Numeric & 
      Lower bin limit for histogram dimension $n$ & \undef{} \\
    $n * 2 + 2$ & Numeric & 
      Upper bin limit for histogram dimension $n$ & \undef{}
  \end{tabular}
\end{center}

The number of histogram dimensions is equal to $(m - 1) / 2$, where
$m$ is the number of SIE dimensions.

Overflow bins are explicitly specified; a negative overflow bin will
have a lower bin limit of $-\infty$, while a positive overflow bin
will have an upper bin limit of $\infty$.

A data sample will fall into a bin if it is in the range
[\emph{lower}, \emph{upper}); the lower bound is closed while the
upper bound is open.

If a bin is specified more than once, the last count is valid.  This
is to allow ``streaming'' histogram data.

The \libsie{} library offers a convenience interface to access a
histogram stored in this schema as an $n$-dimensional array rather
than as a linear list of bins.  See the \libsie{} documentation for
more information.

\subsection{\xml{somat:rainflow} data schema}

This schema is identical to the \xml{somat:histogram} data schema save
for the presence of a metadata tag
\xml{somat:rainflow\_unclosed\_cycles}.  This tag contains a sequence
of ASCII-formatted floating-point numbers containing the unclosed
cycles stack from the rainflow counting algorithm.  Note that the
histogram emitted by SIE has been closed; to open it again the
rainflow algorithm must be run in reverse with the unclosed cycles
stack.

\appendix

\section{License}

\begin{verbatim}
Copyright (C) 2005-2010  HBM Inc., SoMat Products

This document is free software; you can redistribute it and/or
modify it under the terms of version 2.1 of the GNU Lesser General
Public License as published by the Free Software Foundation.

This document is distributed in the hope that it will be useful, but
WITHOUT ANY WARRANTY; without even the implied warranty of
MERCHANTABILITY or FITNESS FOR A PARTICULAR PURPOSE.  See the GNU
Lesser General Public License for more details.

You should have received a copy of the GNU Lesser General Public
License along with this document; if not, write to the Free Software
Foundation, Inc., 51 Franklin Street, Fifth Floor, Boston, MA
02110-1301 USA
\end{verbatim}

\end{document}
